\documentclass[peerreviewcls]{IEEEtran}
\IEEEoverridecommandlockouts
% The preceding line is only needed to identify funding in the first footnote. If that is unneeded, please comment it out.
%Template version as of 6/27/2024


\usepackage{cite}
\usepackage{amsmath,amssymb,amsfonts}
\usepackage{algorithmic}
\usepackage{graphicx}
\usepackage{textcomp}
\usepackage{xcolor}

\usepackage{amsthm}

\usepackage{mathtools}
\usepackage{tikz}
\usepackage{tikz-network}
\usepackage{xspace}
\usepackage{xargs}
\usepackage{fontenc}

\usepackage{subcaption}
\usepackage{caption}

\usepackage{tikz-cd}
\usepackage{tikzit}
% to reduce tikz picture white space
\usepackage{adjustbox}

% for quotienting
\usepackage{xfrac} 


\newtheorem{theorem}{Theorem}[section]
\newtheorem{definition}[theorem]{Definition}
\newtheorem{example}[theorem]{Example}
\newtheorem{proposition}[theorem]{Proposition}
\newtheorem{remark}[theorem]{Remark}

% some macros that maybe should be moved elsewhere

\newcommand\HypI[1]{\textbf{HypI}(#1)}
\newcommand\MdaCospans{\textbf{MHypI}(\Sigma)}
\newcommand{\Ecospans}{{\catname{EHypI(\Sigma)}}}
\newcommand{\MdaEcospans}{{\catname{MEHypI(\Sigma)}}}
\newcommand{\WellTypedMdaEcospans}{{\catname{MEHypI(\Sigma)}}}
% conflict
\newcommand{\hashtag}{{\#}}
\newcommand{\consistency}{{\smile}}

% !TEX root = ../main.tex

\newcommand{\missing}[1]{{\color{red}\bfseries [TODO]}}

\newcommand{\egg}{\texorpdfstring{\MakeLowercase{\texttt{egg}}}{\texttt{egg}}\xspace}
\newcommand{\Egg}{\texorpdfstring{\MakeLowercase{\texttt{egg}}}{\texttt{egg}}\xspace}
% \newcommand{\egg}{Lego\xspace}
% \newcommand{\Egg}{Lego\xspace}
\newcommand{\egraphs}{\mbox{e-graphs}\xspace}
\newcommand{\egraph}{\mbox{e-graph}\xspace}
\newcommand{\Egraph}{\mbox{E-graph}\xspace}
\newcommand{\Egraphs}{\mbox{E-graphs}\xspace}
\newcommand{\eclass}{\mbox{e-class}\xspace}
\newcommand{\Eclass}{\mbox{E-class}\xspace}
\newcommand{\enode}{\mbox{e-node}\xspace}
\newcommand{\eclasses}{\mbox{e-classes}\xspace}
\newcommand{\enodes}{\mbox{e-nodes}\xspace}
\newcommand{\Enodes}{\mbox{E-nodes}\xspace}
% \newcommand{\regraph}{Regraph\xspace}
% \newcommand{\Regraph}{Regraph\xspace}
\newcommand{\sz}{Szalinski\xspace}
\newcommand{\find}{\texttt{find}\xspace}

\newcommand{\equivid}{\equiv_{\sf id}}
\newcommand{\equivnode}{\equiv_{\sf node}}
\newcommand{\equivterm}{\equiv_{\sf term}}

\newcommand{\congrinv}{$\mathcal{I}_c$\xspace}

\newcommand{\egglogo}[1][]{\protect\includegraphics[height=1em, #1]{egg.png} }
\newcommand{\eggurl}{\url{https://github.com/mwillsey/egg}}

% TODOs
% https://tex.stackexchange.com/questions/9796/how-to-add-todo-notes

\newcommand{\eggtodo}[1]{\textcolor{Red}{\textbf{[CITE: #1]}}}

\newcommand{\defTodo}[2]{%
  \expandafter\newcommand\csname #1\endcsname[1]{%
    \todo[linecolor=#2,backgroundcolor=#2!25,bordercolor=#2,inline,size=\tiny]{\textbf{#1}: ##1}}}

\newcommand{\defTODO}[2]{%
  \expandafter\newcommand\csname #1\endcsname[1]{%
    \todo[linecolor=#2,backgroundcolor=#2!25,bordercolor=#2,inline,size=\tiny,caption={\textbf{(#1 LONG TODO)}}]{##1}}}

\newcommand{\AlekseiColor}{RoyalBlue}

% examples
% \newcommand{\RemyColor}{Red}
% \newcommand{\OliverColor}{OliveGreen}
% \newcommand{\ChandraColor}{Magenta}
% \newcommand{\PavelColor}{Cerulean}
% \newcommand{\ZachColor}{Plum}
% \newcommand{\BenColor}{Orchid}
% \newcommand{\JamesColor}{Salmon}


\defTodo{Aleksei}{\AlekseiColor}
% \defTodo{Max}{\MaxColor}

\defTODO{ALEKSEI}{\AlekseiColor}
% \defTODO{MAX}{\MaxColor}


% categorical stuff

\newcommand{\catname}[1]{\mathbf{#1}}

\newcommand{\Graph}{\catname{Graph}}
\newcommand{\Set}{\catname{Set}}

\newcommand\id{\textsf{id}}
\newcommand\sym{\textsf{sym}}

% tiny inline string diagrams
\newcommand{\comonoid}{%
	\begin{tikzpicture}[baseline=0.4ex, scale=0.4, line width=0.8pt]
		\path [use as bounding box] (0,0) rectangle (0.6,0.8);
		\draw (0.2,0) -- (0.2,0.4);
		\filldraw[black] (0.2,0.4) circle (0.08);
		\draw (0,0.8) .. controls (0,0.65) and (0.1,0.45) .. (0.2,0.4);
		\draw (0.4,0.8) .. controls (0.4,0.65) and (0.3,0.45) .. (0.2,0.4);
	\end{tikzpicture}%
}
\newcommand{\counit}{%
	\begin{tikzpicture}[baseline=0.4ex, scale=0.4, line width=0.8pt]
		\path [use as bounding box] (0,0) rectangle (0.4,0.8);
		\draw (0.2,0) -- (0.2,0.4);
		\filldraw[black] (0.2,0.4) circle (0.08);
	\end{tikzpicture}%
}

% camera-ready
\newcommand{\ifcameraready}[2]{\ifdefined\cameraready #2 \else #1 \fi}

\input{../../sample.tikzstyles}
\input{../../hypergraph.tikzstyles}
\input{../../hypergraph.tikzdefs}


\def\BibTeX{{\rm B\kern-.05em{\sc i\kern-.025em b}\kern-.08em
    T\kern-.1667em\lower.7ex\hbox{E}\kern-.125emX}}
\begin{document}

\title{E-graphs With Bindings
}

% \author{\IEEEauthorblockN{1\textsuperscript{st} Given Name Surname}
% \IEEEauthorblockA{\textit{dept. name of organization (of Aff.)} \\
% \textit{name of organization (of Aff.)}\\
% City, Country \\
% email address or ORCID}
% \and
% }

\author{Anonymous}

\maketitle

\begin{abstract}
This is abstract.
\end{abstract}

\begin{IEEEkeywords}
component, formatting, style, styling, insert.
\end{IEEEkeywords}

\section{Introduction}

\subsection{Motivation and contribution}

E-graphs (short for \emph{equality graphs})~\cite{EggPaper} are a data structure incorporating equivalence representation and used in automated theorem proving, compiler optimisation, and program analysis. 
Using e-graphs mitigates the so-called \emph{phase-ordering problem} of term rewriting, that is the fact that the way rewrites are scheduled may enable or block further rewrites due to introduction or elimination or redexes. 
This is achieved by making rewrites on e-graphs non-destructive: an application of a rewrite rule instead of replacing the left-hand side by the right-hand side, adds the latter to the graph so that all pre-existing redexes are protected and new ones are introduced. 
Explosive growth of the graph is avoided by the clever use of sharing of vertices. 

E-graphs were first introduced in the automated theorem proving literature in the 1980s~\cite{nelson1980techniques}, but recent years have seen renewed interest particularly for optimisation and synthesis using so called \emph{equality saturation}~\cite{10.1145/1594834.1480915, griggio_proceedings_2022, EggPaper,flatt_small_2022}, the repeated application of rewrites to an e-graph until a fixpoint is reached. 
The wide range of application domains led to extensions of the e-graph formalism, for instance to account for conditional rewriting~\cite{singher2023colored} or  to combine e-graphs with database queries (\emph{relational e-matching})~\cite{zhang_relational_2022}.
Of particular interest to us is using e-graphs in the context of functional programming languages, which requires a representation of the $\lambda$-calculus using it. 
Pioneering work in this direction has been done by~\cite{koehler2022sketchguided}.
The well-known twin challenges of binding and capture-avoiding substitution remain however somewhat problematic in this work because DAN: IT IS UNCLEAR WHY THIS IS PROBLEMATIC IN LOC. CIT. THE TEXT BELOW IS UNCONVINCING. CAN YOU ILLUSTRATE THE PROBLEM WITH SOME SIMPLE EXAMPLES? 
\begin{quotation}
There are domains, however, where the effective usage of e-graphs requires \textit{ad hoc} tweaks to achieve the desired performance, namely domains that include variable \textit{binding} and the closely related notion of \textit{substitution}, for example, lambda calculus with $\beta$-reduction.
The first issue appears due to the fact that lambda calculus terms are quotiented by renaming of bound variables.
If such renaming is not accounted for by the language encoding, \textit{e.g.,} by using \textit{de Brujin} indices, populating an e-graph with $\alpha$-equivalent terms increases the size of an e-graph without adding any meaningful equivalences.
The other issue is that having binding entails having a substitution and computing substitution in e-graphs is hard as it is not a mere substitution of terms but rather a substitution of equivalence classes that represent a lot of terms at once.
\end{quotation}

In this paper we take what we consider to be a principled approach to e-graphs with bindings by lifting to monoidal closed categories prior work showing how e-graphs (without binding) can be given a categorical semantics using semilattice-enriched symmetric monoidal categories, which allows e-graphs to be reconstructed as string diagrams.~\cite{ghica2024equivalencehypergraphsegraphsmonoidal}  
The approach in \emph{loc. cit.} was originally motivated by the desire to support the machinery of categorial (DPO) graph rewriting for e-graphs, but by extending the semilattice-enrichment to monoidal-categories it paves the way to a natural derivation of e-graphs with binding. 

\subsection{String diagrams}

\subsection{E-graphs}

E-graphs are a data structure which can efficiently represent many equivalent terms of an algebraic theory.  They generalise the typical DAG representation of terms to include equivalence classes.
The key concepts of e-graphs can be intuitively grasped from some simple (but non-trivial) examples, as shown in Figure~\ref{fig:e-graph-example}.
This example is a `Hello world!' of e-graphs, presented first in~\cite{EggPaper}.

\begin{figure*}
\[
\adjustbox{scale=0.7}{
\tikzfig{figures/e-graph-example}
}
\]
\caption{E-graph example (top) and its equivalent string diagram representation (bottom)}
\label{fig:e-graph-example}
\end{figure*}

Each column represents a term (or term rewrite rule), the conventional e-graph representation, and the equivalent string diagram representation, an e-hypergraph as developed in~\cite{ghica2024equivalencehypergraphsegraphsmonoidal}. 
The initial term, $(a*2)/2$ (which corresponds to subfigure (a)), is already represented efficiently in the e-graph by sharing the node 2.
The string diagram version is slightly different than conventional e-graphs by representing the sharing of the node 2 explicitly, using a sharing (contraction) node
\scalebox{0.3}{
	\begin{tikzpicture}
		\begin{pgfonlayer}{nodelayer}
			\node [style=vertex] (0) at (1.5, 3.5) {};
			\node [style=none] (1) at (1.5, 3) {};
			\node [style=none] (2) at (1, 4) {};
			\node [style=none] (3) at (2, 4) {};
		\end{pgfonlayer}
		\begin{pgfonlayer}{edgelayer}
			\draw (0) to (1.center);
			\draw [bend right=45] (0) to (3.center);
			\draw [bend left=45] (0) to (2.center);
		\end{pgfonlayer}
	\end{tikzpicture}
}
. 

Nodes of an e-graph are encapsulated in dashed boxes indicating equivalence classes of nodes, or, more generally, of subgraphs rooted at these boxes.
Initially each equivalence class has a single node (subgraph).
The first rewrite rule, replacing multiplication by the more efficient \emph{shift-left} operator creates the first non-unitary (containing more than one node or subgraph) equivalence class, which includes the multiplication ($*$) and shift-left ($<\!\!<$) operator nodes, with a new sharing, that of node $a$.
This second e-graph illustrates the other difference, besides sharing, between conventional e-graphs and e-hypergraphs, namely the fact that in the former edges connect directly to nodes inside the equivalence class, whereas in the latter edges connect to the equivalence class itself.
Explicit discarding (weakening) nodes, depicted with black dots with a single wire inside the dashed boxes 
\scalebox{0.4}{
	\begin{tikzpicture}
		\begin{pgfonlayer}{nodelayer}
			\node [style=vertex] (0) at (1.5, 3.5) {};
			\node [style=none] (1) at (1.5, 3) {};
		\end{pgfonlayer}
		\begin{pgfonlayer}{edgelayer}
			\draw (0) to (1.center);
		\end{pgfonlayer}
	\end{tikzpicture}
}, indicate which of the class-level edges are connected to which nodes inside the class.

The other steps, (c) to (e), represent further elaborations of the e-graph, or e-hypergraph, via the application of more rules.


A typical way on encoding $\beta$-reduction in lambda-calculus in e-graphs is by explicit substitution~\cite{EggPaper,koehler2022sketchguided}: by adding a special substitution nodes and introduce rewrite rules for them.
After the reduction has been done these nodes serve no purpose and can be removed.
We show that $\beta$-reduction comes naturally as a composition when using string diagrams for enriched closed monoidal categories.
Consider an example in Figure~\ref{fig:e-graph-substitution} where successive iterations of an e-graph for a term $(\lambda x . (y + y) + x) 1$ are shown following the $\beta$-reduction.
First the subterm for an explicit substitution is introduced which is deemed equivalent to the top-level application term.
It is then rewritten using the rules for explicit substitution at the top of each column.
At the end all explicit substitution nodes can be removed.

Now consider an example from Figure~\ref{fig:e-string-substitution} which encodes the same $\beta$-reduction for the same term.
The $\lambda$-abstraction is represented by a rounded box and the bound variable is the dangling wire.
The application node is depicted as 
\adjustbox{scale=0.15}{\begin{tikzpicture}
	\begin{pgfonlayer}{nodelayer}
		\node [style=none] (0) at (1, 6.75) {};
		\node [style=none] (1) at (1.25, 7.75) {};
		\node [style=none] (2) at (2, 8) {};
		\node [style=none] (3) at (2.75, 7.75) {};
		\node [style=none] (4) at (3, 6.75) {};
		\node [style=none] (5) at (1.5, 6.75) {};
		\node [style=none] (6) at (2.5, 6.75) {};
		\node [style=none] (7) at (2, 9) {};
		\node [style=none] (8) at (1.5, 5.75) {};
		\node [style=none] (9) at (2.5, 5.75) {};
	\end{pgfonlayer}
	\begin{pgfonlayer}{edgelayer}
		\draw [style=new edge style 2] (3.center)
			 to [in=60, out=-30, looseness=0.75] (4.center)
			 to (0.center)
			 to [in=-150, out=120, looseness=0.75] (1.center)
			 to [in=-180, out=45, looseness=0.75] (2.center)
			 to [in=135, out=0, looseness=0.75] cycle;
		\draw (2.center) to (7.center);
		\draw (5.center) to (8.center);
		\draw (9.center) to (6.center);
	\end{pgfonlayer}
\end{tikzpicture}
}. 
The reduction is just the re-attachment of wires.
Note how we avoided re-attachment for subterms as they did not use the bound variable wire, unlike in the case $[x \mapsto 1](y+y)$ before.

A common solution to replace explicit variable names in e-graphs is by using De Brujin indices~\cite{koehler2022sketchguided}.
The issue with this approach is that it interferes with sharing as equivalent subterms can be structurally different.
Consider the example $\lambda f . f ((\lambda x . f x) 2)$ which is represented as $\lambda . \%0 ((\lambda . \%1 \%0) 2)$ using De Brujin indices.
Notice how two different occurrences of the same bound variable $f$ are represented differently by $\%0$ and $\%1$ which means they can not be shared in an e-graph.
A string diagrammatic representation of the same term makes it possible for $f$ to be shared as demonstrated below.

\[
\adjustbox{scale=0.6}{
    \tikzfig{figures/de-brujin-string}
}
\]

Using De Brujin indices when performing $\beta$-reduction also introduces the complexity of representing the reduction as a rewrite rule.
Special shifting operators need to be introduced on the right-hand side of a rewrite rule which makes it not merely syntactic~\cite{koehler2022sketchguided}.


\textcolor{red}{Add example showing that using De Brujin indices with e-graphs breaks sharing}

\textcolor{red}{Add example showing $\beta$-reduction in e-graphs compared to string diagrams}


\begin{figure*}
    \[
        % \hspace{1.3cm}
        \adjustbox{width=\textwidth}{
        \tikzfig{figures/e-graph-substitution}
        }
    \]
    \caption{E-graph explicit substitution example.}
    \label{fig:e-graph-substitution}
  \end{figure*}
  
  \begin{figure*}
    \[
        % \hspace{1.3cm}
        \adjustbox{width=\textwidth}{
        \tikzfig{figures/e-string-substitution-2}
        }
    \]
    \caption{String diagrammatic substitution example.}
    \label{fig:e-string-substitution}
  \end{figure*}

% !TEX root = ../main.tex

\newtheorem*{notation*}{Notation}

\section{Categorical Semantics}

\subsection{Symmetric Monoidal Categories and String Diagrams}

We will assume familiarity with the definition of category,  functor and natural transformation.
\begin{notation*}
We use the following notation: given a catgeory $\mathbb{C}$, we write $ob(\mathbb{C})$ for the set of objects $\mathbb{C}(A,B)$ for hom-sets and $f: A \to B$ for $f \in \mathbb{C}(A,B)$. We write $\textbf{id}_A$ for the identity on $A$ and commonly write $f;g$ for composition in diagram order.  Composition in the usual order is written $g \circ f$. 
We will often omit subscripts on identities and natural transformations where it can be inferred. 
\end{notation*}
\begin{definition}
A \textit{strict symmetric monoidal category (SMC)} $(\mathbb{C}, \otimes, 1)$ is given by a category $\mathbb{C}$ with the following equipment: 
\begin{itemize}
\item a \textit{tensor product} functor $\otimes: \mathbb{C} \times \mathbb{C} \to \mathbb{C}$,
\item a \textit{tensor unit} object $1 \in ob(\mathbb{C})$, 
\item the \textit{symmetry} natural isomorphism: $\sigma: A \otimes B \to B \otimes A$
\item and such that: $A \otimes (B \otimes C) = (A \otimes B) \otimes C$ ,  $1 \otimes A = A$ ,  $A \otimes 1 = A$
\end{itemize}
satisfying certain coherence conditions [MacLane].  A \textit{permutations and products category (PROP)} is a symmetric monoidal category where every object is the n-fold tensor product of some distinguished object $A$. In this case, we will enote the n-fold tensor product of the distinguished object by simply $n$, so that $n \otimes m = n+m$.

A \textit{strict symmetric monoidal functor} $F: (\mathbb{C}, \otimes_\mathbb{C}, 1_\mathbb{C}) \to (\mathbb{D}, \otimes_\mathbb{D}, 1_\mathbb{D})$ is given by an underlying functor $F: \mathbb{C} \to \mathbb{D}$ such that 
\[
 1_{\mathbb{D}} = F(1_{\mathbb{C}})
 \qquad F(A \otimes_{\mathbb{C}} B) = F(A) \otimes_{\mathbb{D}} F(B)\ ,
\]
satisfying certain coherence conditions. 
We adopt the convention that $\otimes$ binds more tightly than $(;\!)$. 
A \textit{PROP functor} is additionally identity-on-objects. 
\end{definition}

String diagrams for {SMCs} are constructed as in Figure 1. Note, $\textsf{id}_0$ is denoted by the empty diagram. 


\subsection{Symmetric Monoidal Theories and PROPs. }

A presentation of an algebraic theory is traditionally given by a set together with a \textit{signature} of $n$-ary operations and a set of \textit{equations} -- formally,  pairs of terms freely generated over the signature.  

For example, the \textit{semilattice signature} consists of a single binary operation $+$.  Given a set $A$, this can be used to freely generate terms according to the following grammar.
\begin{align*}
	\infer{x_1, \ldots, x_n \vdash x_j}{x_i \in A} \qquad \infer{\Gamma \vdash t_1 + t_2}{\Gamma \vdash t_1 && \Gamma \vdash t_2} 
\end{align*}
Such terms may be organized into a Cartesian category, with morphisms given by tuples of terms of the following form,  and composition is given by substitution of terms for free variables.  
\[
	x_1, \ldots, x_n \vdash (t_1, \ldots, t_m) : n \to m\ .
\]
A set of equations can then be given as a set of parallel pairs of morphisms.  
For example, the equations for the theory of semilattices are the following.
\[
	x,y,z \vdash x+(y+z) = (x+y)+z : 3 \to 1\ \qquad  
    x,y \vdash x+y = y+x \qquad 
    x \vdash x+x = x
\]
Now, quotienting the category by the least congruence identifying every such parallel pair, we recover a category whose morphisms are \textit{equivalence classes} of terms modulo the expected equations. 
This gives a categorical presentation of an algebraic theory, known as its \textit{Lawvere category}. 
Importantly,  its category of models (\textit{i.e.}, category of functors to \textbf{Set}) is equivalent to the expected category of models of the algebraic theory in question.  

Presenting algebraic theories as Lawvere categories gives a framework in which the concept is amenable to generalization.  We are interested here in symmetric monoidal theories, which are generalizations of algebraic theories where the operations of the signature may have arbitrary \textit{co-arities}.  Lawvere categories will be generalized to particular symmetric monoidal categories called PROPs, but the ideas are otherwise as described above. 
\begin{definition}
A \textit{monoidal signature} $\Sigma$ is given by a set of \textit{generators} $f: n \to m$,  with \textit{arity} $n$ and \textit{coarity} $m$.  A \textit{signature homomorphism} $h: \Sigma \to \Gamma$ is a function from $\Sigma$ to $\Gamma$ which preserves the (co-)arities of elements. 
\end{definition}
Analogous to freely generating terms over an algebraic signature,  we now consider terms freely generated over a monoidal signature.  These are naturally arranged into a particular form of symmetric monoidal category: a PROP. 

We can form a category $\textbf{Sig}$ of monoidal signatures and signature homomorphisms.  There is an evident forgetful functor $U: \textsf{PROP} \to \textsf{Sig}$ from the category of PROPs, which considers every morphism of a PROP as an operation of a signature, forgetting the remaining structure of the category.   Its left adjoint $F$, pictured below, left,  generates a free PROP over a monoidal signature.
% https://q.uiver.app/#q=WzAsNixbMCwxLCJcXHRleHRiZntTaWd9Il0sWzIsMSwiXFx0ZXh0YmZ7UFJPUH0iXSxbMSwxLCJcXHJvdGF0ZWJveHs5MH17XFx2ZGFzaH0iXSxbMywwLCJcXFNpZ21hIl0sWzUsMCwiVUYoXFxTaWdtYSkiXSxbNSwyLCJVKFxcbWF0aGJie0N9KSJdLFswLDEsIkYiLDAseyJjdXJ2ZSI6LTJ9XSxbMSwwLCJVIiwwLHsiY3VydmUiOi0yfV0sWzMsNCwiaSIsMCx7InN0eWxlIjp7InRhaWwiOnsibmFtZSI6Imhvb2siLCJzaWRlIjoidG9wIn19fV0sWzMsNSwiXFxmb3JhbGwgaiIsMix7InN0eWxlIjp7InRhaWwiOnsibmFtZSI6Imhvb2siLCJzaWRlIjoidG9wIn19fV0sWzQsNSwiXFxleGlzdHMgVShcXGxsYnJhY2tldCAtIFxccnJicmFja2V0KSJdXQ==
\[\begin{tikzcd}
	&&& \Sigma && {UF(\Sigma)} \\
	{\textsf{Sig}} & {\ \  \perp} & {\textsf{PROP}} \\
	&&&&& {U(\mathbb{C})}
	\arrow["F", bend left, from=2-1, to=2-3]
	\arrow["U", bend left, from=2-3, to=2-1]
	\arrow["i", hook, from=1-4, to=1-6]
	\arrow["{\forall j}", from=1-4, to=3-6]
	\arrow["{\exists U(\llbracket - \rrbracket)}", from=1-6, to=3-6]
\end{tikzcd}\]
The corresponding universal property is given above, right.  The signature morphism $i:\Sigma \to UF(\Sigma)$ is given by the obvious inclusion.  Roughly,  for every PROP $\mathbb{C}$, a signature homomorphism $f: \Sigma \to U(\mathbb{C})$ , which chooses a term of $\mathbb{C}$ for each generator $\Sigma$,  uniquely determines a functor $\llbracket - \rrbracket: F(\Sigma) \to \mathbb{C}$ which preserves this choice.  In the sequel, we will name the free functor $F$ as $\textsf{PROP}(-): \textsf{Sig} \to \textsf{PROP}$. 

\begin{definition}
The free \textit{products and permutations category} over a monoidal signature $\Sigma$ is $\textsf{PROP}(\Sigma)$.  
\end{definition}
The free category $\textsf{PROP}(\Sigma)$ can be syntactically constructed from taking (typed) categorical combinator terms freely constructed from generators $f \in \Sigma$, $\textsf{id}_1$, $\sigma_{1,1}$, $(;\!)$ and $\otimes$, quotiented by the axioms of a symmetric monoidal category.  

We can now specify an equation associated with a monoidal signature as a pair of parallel morphisms of the form above, leading to the following definition.
\begin{definition}
A \textit{presentation of a symmetric monoidal theory} $(\Sigma, \mathcal{E})$ is given by a pair of a monoidal signature $\Sigma$ and a set $\mathcal{E}$ of pairs of parallel morphisms $\phi,\psi: n \to m$. 
\end{definition}
Finally, analogous to a Lawvere theory, we define a \textit{symmetric monoidal theory}. 
\begin{definition}
A \textit{symmetric monoidal theory} $\textsf{SMT}(\Sigma,\mathcal{E})$ generated by a presentation $(\Sigma, \mathcal{E})$ is given by $\textsf{PROP}(\Sigma)$ quotiented by the least congruence including $\mathcal{E}$.
\end{definition}

It will often be convenient to present SMTs in terms of string diagrams [which are technically equiv classes of arrows?!] 
\begin{example}
The SMT of \textit{commutative comonoids} is given by the following generators, ${\Delta, !}$, depicted in string diagrammatic form:
\[
    \tikzfig{categorical-semantics/Cartesian-equipment}
\]
together with the following associativity, commutativity and unitality equations, again given in terms of string diagrams. 
\[
    \tikzfig{categorical-semantics/Cartesian-theory}
\]
\end{example}

\begin{example}
The Cartesian SMT is given by a set of generators $\Sigma_C$, together with the generators and equations of the SMT of commutative monoids, and the following additional naturality equations, for every $c \in \Sigma_C$
\[
    \tikzfig{categorical-semantics/Cartesian-naturality}
\]
By Fox's Theoem [cite],  $\textsf{SMT}(\Sigma_C, \mathcal{E}_C)$ is equivalent to the free Cartesian category over $\Sigma_C$. 
\end{example}

\subsection{Semilattice Enriched Monoidal Categories}

We will use enrichment to axiomatize the e-box structure as the ``join" or``sum" of two morphisms $f,g: A \to B$ as $f + g: A \to B$.  More precisely our hom-sets will be given the structure of a semilattice.  

Generally,  we may define a category \textit{enriched} in any other monoidal category $M$ (\textit{e.g.},  \textbf{Semilattice}).  In this case,  hom-sets are generalized to hom-objects of $M$ (\textit{e.g.},  a semilattice) and composition 
\[
	\circ: Hom(B,C) \otimes Hom(A,B) \to Hom(A,C)
\]
is defined as a $M$-morphism.  In other words,  composition must respect the additional structure on hom-sets.  
In our case,  the category we wish to enrich is also monoidal,  so we additionally ask that the monoidal structure also respects this structure.  Technically, this amounts to asking that the tensor product is an enriched functor. 
We omit further details of the general case of enrichment here, and instead work with the following more concrete defintion, where $M = \textbf{Semilattice}$. 


%\begin{definition}
%A \textit{semigroup} is a pair $(A, +)$ of a set $A$ together with an associative binary operation $+: A \times A \to A$.  We say a semigroup is \textit{commutative} when $+$ is {commutative} and \textit{idempotent} when $+$ is idempotent.  
%Because we will only be concerned with commutative, idempotent semigroups in this paper, we will now often abbreviate the former to simply ``semigroup". 
%\end{definition}

\begin{definition}
A \textit{semilattice enriched strict SMC (PROP)}  $\mathbb{C}$ is given by an SMC (PROP) $(\mathbb{C}, \otimes, 1,+)$ where every hom-set additionally has the structure of a semilattice $(\mathbb{C}(A,B), +)$ which respects the composition and tensor product in the following ways,  for all appropriately typed morphisms. 
\begin{align*}
f ; (g+h) &= f;g + f;h &
(f+g) ; h &= f;h + g;h \\
f \otimes (g+h) &= f \otimes g + f \otimes h & 
(f+g) \otimes h &= f \otimes h + g \otimes h,
%\item A set $ob(\mathbb{C})$ of objects;
%\item For every pair $A,B \in ob(\mathbb{C})$, a hom-semigroup $Hom(A,B)$;
%\item An element $\textsf{id}_{A,B} \in Hom(A,B)$; 
\end{align*}
Note,  we take $(;\!)$ and $\otimes$ to bind more tightly than $+$.

A \textit{semilattice enriched strict SMC (PROP) functor} $F: (\mathbb{C}, \otimes, 1,+) \to (\mathbb{D}, \otimes, 1,+)$ is given by a strict SMC (PROP) functor on the underlying SMCs (PROPs) which additionally satisfies $F(f+g) = F(f)+F(g)$. 
\end{definition}
Note, in particular,  that semigroups \update{semilattice?} do not require a unit for $+$. {This is important, since, in any commutative monoid enriched category,  if there exists a categorical product, then it must also be a categorical coproduct -- a degeneracy we wish to avoid.}


As before, there is an evident forgetful functor $U: \textsf{PROP}^+ \to \textsf{PROP}$ from the category of semilattice enriched PROPs, which forgets the semilattice structure on hom-sets. Its left adjoint $F$, pictured below, generates the free enrichmenet of a PROP. 
% https://q.uiver.app/#q=WzAsNixbMCwxLCJcXHRleHRzZntQUk9QfSJdLFsyLDEsIlxcdGV4dHNme1BST1B9XisiXSxbMSwxLCJcXHBlcnAiXSxbMywwLCJcXFNpZ21hIl0sWzUsMCwiVUYoXFxTaWdtYSkiXSxbNSwyLCJVKFxcbWF0aGJie0N9KSJdLFswLDEsIkYiLDAseyJjdXJ2ZSI6LTJ9XSxbMSwwLCJVIiwwLHsiY3VydmUiOi0yfV0sWzMsNCwiaSIsMCx7InN0eWxlIjp7InRhaWwiOnsibmFtZSI6Imhvb2siLCJzaWRlIjoidG9wIn19fV0sWzMsNSwiXFxmb3JhbGwgaiIsMix7InN0eWxlIjp7InRhaWwiOnsibmFtZSI6Imhvb2siLCJzaWRlIjoidG9wIn19fV0sWzQsNSwiXFxleGlzdHMgVShcXGxsYnJhY2tldCAtIFxccnJicmFja2V0KSJdXQ==
\[\begin{tikzcd}
	&&& \mathbb{C} && {UF(\mathbb{C})} \\
	{\textsf{PROP}} & \perp & {\textsf{PROP}^+} \\
	&&&&& {U(\mathbb{D})}
	\arrow["F", bend left, from=2-1, to=2-3]
	\arrow["U", bend left, from=2-3, to=2-1]
	\arrow["i", hook, from=1-4, to=1-6]
	\arrow["{\forall j}"', hook, from=1-4, to=3-6]
	\arrow["{\exists U(\llbracket - \rrbracket)}", from=1-6, to=3-6]
\end{tikzcd}
\]
The corresponding universal property is given above, right. The symmetric monoidal functor $i: \mathbb{C} \to UF(\mathbb{C})$ is given by the obvious inclusion of symmetric monoidal categories.  In the sequel, we will name the free functor $F$ as $(-)^+: \textsf{PROP} \to \textsf{PROP}^+$.  We further name the composition of the functors $\textsf{PROP}(-)$ and $\textsf{SMT}(-)$ with $(-)^+$ as $\textsf{PROP}^+(-), \textsf{SMT}^+(-): \textsf{Sig} \to \textsf{PROP}^+$: these generate the free enriched PROP and freely enriched SMT over a monoidal signature and presentation of a symmetric monidal theory, respectively.  

\begin{definition}
The free enriched PROP over a monoidal signature $\Sigma$ is $\textsf{PROP}^+(\Sigma)$. 
The freely enriched SMT over a monoidal signature $\Sigma$ is $\textsf{SMT}^+(\Sigma, \mathcal{E})$. 
\end{definition}
The free category $\textsf{PROP}^+(\Sigma)$ can be syntactically constructed from taking (typed) categorical combinator terms freely constructed from generators $f \in \Sigma$, $\textsf{id}_1$, $\sigma_{1,1}$, $(;\!)$ and $\otimes$, as well as $f+g$, quotiented by the axioms of an enriched symmetric monoidal category.  

To aid reasoning, we introduce a new language of string diagrams for semilattice enriched SMCs, using a hierarchical ``box'' structure to capture the binary operation on morphisms.  This will later be used in the translation of equivalence classes from the e-graph to the string diagrammatic setting. 

Figure \ref{fig:egraph-strings} displays the generators of this language and Figure \ref{fig:string-equations} displays the additional equations which these diagrams satisfy, in addition to the standard SMC equations. 

We use string diagrams informally in this paper, but it is intuitively clear that these diagrams are sound and complete with respect to their intended categorical semantics. Indeed, the diagrammatic language itself is not really novel and is similar to \textit{e.g.}, \cite{duncan_generalised_2009}. 

\begin{figure}
\[  
    \scalebox{0.75}{
	\tikzfig{categorical-semantics/egraph-strings-equations}
    }
\]
\caption{Equations for a  semilattice enriched symmetric monoidal category}
\label{fig:string-equations}
\end{figure}


\section{Combinatorial semantics}


We can extend the definition of e-hypergraphs for $\catname{SLat}$-enriched symmetric monoidal categories~\cite{ghica2024equivalencehypergraphsegraphsmonoidal} to accommodate new structure brought by closedness.

\begin{definition}
Closed monoidal signature $\Sigma$ is a pair $(\Sigma_{O},\Sigma_{M})$ of signatures for objects and signatures for morphisms with sources and targets in $\text{obj}_{\Sigma_{O}}$ defined as
\begin{itemize}
  \item $I$ (monoidal unit) is in $\text{obj}_{\Sigma_{O}}$
  \item For every $o \in \Sigma_{O}$, $o \in \text{obj}_{\Sigma_{O}}$;
  \item For every $o_1$ and $o_2$ $ \in \text{obj}_{\Sigma_{O}}$, $o_1 \otimes o_2 \in \text{obj}{\Sigma_{O}}$ and $o_1 \multimap o_2 \in \text{obj}_{\Sigma_{O}}$
\end{itemize}
\end{definition}

\begin{definition}
 Closed symmetric monoidal theory ($\catname{CSMT}$) over a closed monoidal signature $\Sigma$ and equations $\mathcal{E}$ denoted as $\catname{CSMT}(\Sigma, \mathcal{E})$ is given by a set of terms $T$
 \begin{align*}
 T \coloneq\; &t : A \to B \in \Sigma_{M} \;|\; t_1 \otimes t_2 : A \otimes C \to B \otimes D \text{ for } t_{1} : A \to B, t_{2} : C \to D \in T\\
            & \;|\; t_{1};t_{2} : A \to C \text{ for } t_{1} : A \to B, t_{2} : B \to C \in T \; | \;  \Lambda(t) : A \to X \multimap B \text{ for } t : A \otimes X \to B \\
            &  \;|\; \text{eval}(t_1 \otimes t_2) : B \text{ for } t_1 : X \multimap B, t_2 : X \in T \\
            &  \;|\; \text{sym}_{A,B} \;|\; id_{A}
 \end{align*}
 quotiented by the laws of a closed symmetric monoidal category.
\end{definition}
% \question{Maybe types and terms combined is not a good idea}
\begin{definition}

We define an e-hypergraph $\mathcal{G}$ over a closed symmetric monoidal signature $\Sigma$ as a tuple $(V,E,s,t, l_{V}, l_{E} \textcolor{red}{<}, \textcolor{blue}{<},\consistency)$ where
\begin{itemize}
  \item $V$ is a set of vertices;
  \item $E = \textcolor{gray}{E} \cup \textcolor{red}{E} \cup \textcolor{blue}{E}$ is a set of hyperedges that include \textit{plain} edges and two types of \textit{hierarchical} edges respectively;
  \item $s,t : E \to V^{*}$ are source and target functions;
  \item $l_{E} : E \to \Sigma_{M} + 1$, $l_{V} : V \to \text{obj}_{\Sigma_{O}}$ are \textit{label} functions;
  \item $\textcolor{red}{<} : V + E \to \textcolor{red}{E}$ and $\textcolor{blue}{<} : V + E \to \textcolor{blue}{E}$ are partial \textit{child} functions that are disjoint: if $\textcolor{red}{<}(x)$ is defined then $\textcolor{blue}{<}(x)$ is undefined and vice-versa;
  \item $\consistency = \bigcup_{p} \consistency_{p}$, where $\consistency_{p}$ is a \textit{consistency} relation defined on each set $\{x \in V + E \;|\; p \;\textcolor{red}{<^{\mu}}\; x\}$;
\end{itemize} 
\end{definition}

$\textcolor{red}{E}$ will contain edges that encode the equivalence classes and $\textcolor{blue}{E}$ the ones that defined the lambda abstraction.
All the above functions and relations must satisfy the obvious properties (TODO: add them).

Similarly, we have a category of (extended) cospans formed by such e-hypergraphs where morphisms are of the form
\[
n \xrightarrow{f_{\text{ext}}} n' \xrightarrow{f_{\text{int}}} \mathcal{G} \xleftarrow{g'_{\text{int}}} m' \xleftarrow{g_{\text{ext}}} m
\]

where $n$ and $m$ are external input and output interfaces respectively and $n'$, $m'$ are input and output internal interfaces.
$n,n',m,m'$ are discrete ordered e-hypergraphs.
We further require that $f_{\text{ext}}$ and $g_{\text{ext}}$ are monos and all vertices in the image of $f_{\text{ext}};f_{\text{int}}$ (respectively, $g_{\text{ext}};g_{\text{int}}$) are top-level.

To build a correspondence between terms of a closed SMC and e-hypergraphs with interfaces (above) we restrict the cospans of the latter to monogamous cospans.

\begin{definition}
  We call a cospan 
  \[
n \xrightarrow{f_{\text{ext}}} n' \xrightarrow{f_{\text{int}}} \mathcal{G} \xleftarrow{g'_{\text{int}}} m' \xleftarrow{g_{\text{ext}}} m
\]
\textit{weak} monogamous if
\begin{itemize}
  \item in- and out- degrees of every vertex is at most 1;
  \item $f_{\text{int}}$ and $g_{\text{int}}$ are monos;
  \item Vertices with in-degree (respectively, out-degree) of 0 are precisely the image of $f_{\text{int}}$ (respectively, $g_{\text{int}}$).
\end{itemize}

\end{definition}


\begin{definition}[Isomorphic cospans]
Consider a relation 
\[
R = \{ x R y \text{ if } f_{\text{int}}(x) \consistency f_{\text{int}}(y) \;\text{or}\; \exists z \;\text{s.t.}\; z \textcolor{blue}{<^{\mu}} f_{\text{int}}(x) \text{ and } z \textcolor{blue}{<}^{\mu} f_{\text{int}}(y)\}
\]
for $x, y \in n'$.
And let $S$ be its reflexive closure.
The latter partitions $n'$ into non-empty subsets $\{p_{j}\}^{k}_{j=1}$.
We get an analogous partition for $m'$.

Two extended cospans are isomorphic if there exist isomorphisms $\alpha$, $\beta$ and $\gamma$ making the following diagram commute and such that $\alpha$ and $\gamma$ preserve order within each $p_j$.
\[
\scalebox{0.8}{
    \tikzfig{../../figures/combinatorial_semantics/isomorphic_e_cospans}
}
\]
\end{definition}


Consider an example of isomorphic and non-isomorphic cospans below and their respective string diagrams

\begin{figure}
  \begin{subfigure}[c]{0.4\linewidth}
    \[
    \adjustbox{scale=0.6}{
    \tikzfig{./figures/closed_iso_example_1}
    }
    \]
  \end{subfigure}
  \hfill
  \begin{subfigure}[c]{0.4\linewidth}
    \[
    \adjustbox{scale=0.6}{
    \tikzfig{./figures/closed_iso_example_2}
    }
    \]
  \end{subfigure}
\end{figure}

\begin{definition}

We define ordered sets of \textit{input} (respectively, \textit{output}) vertices of a hierarchical edge $e$ as the vertices $v_{i}$ such that $e < v_{i}$ which are in the image of $f_{\text{int}}$ (respectively, $g_{\text{ext}}$) with the induced ordering.
\end{definition}

\begin{definition}

Given an ordered set $S$ of labelled vertices we can form words formed by concatenating the labels.
We will denote such words as $w(S)$.
\end{definition}

\begin{definition}
We call a monogamous cospan
\[
  n \xrightarrow{f_{\text{ext}}} n' \xrightarrow{f_{\text{int}}} \mathcal{G} \xleftarrow{g'_{\text{int}}} m' \xleftarrow{g_{\text{ext}}} m
\]
\textit{well-typed} if all hierarchical edges of $\mathcal{G}$ are well-typed in the sense below.
\begin{itemize}
  \item For each $e \in \textcolor{red}{E}$ consider sets $I$ and $O$ of its input and output vertices partitioned according to $\consistency_{e}$.
        $e$ is well-typed if for each element $S$ of the partition of $I$ $w(S) = w(s(e))$ and similarly for $O$ and $t(e)$.
  \item For each $e \in \textcolor{blue}{E}$ consider sets $I$ and $O$ of its input and output vertices.
        $e$ is well-typed if there exists an object $B \in \text{obj}_{\Sigma_{O}}$ such that $[...s(e), B] = w(I)$ and $B \multimap w(O) = w(t(e))$
\end{itemize}
\end{definition}



\bibliographystyle{acm}
\bibliography{../../bibliography}

\end{document}
